\documentclass{article}
\usepackage{amsmath, amssymb, amsthm}
\usepackage{hyperref}
\usepackage{geometry}
\geometry{margin=1in}

\title{Erd\H{o}s Problem \#436}
\date{Last updated: 2025-08-31}
\author{Source: erdosproblems.com}

\begin{document}
\maketitle

\noindent\textbf{Status:} OPEN \quad \textbf{No prize}

\noindent\textbf{Tags:} number theory

\medskip
\noindent\textbf{Problem Statement:}

If $p$ is a prime and $k,m\geq 2$ then let $r(k,m,p)$ be the minimal $r$ such that $r,r+1,\ldots,r+m-1$ are all $k$th power residues modulo $p$. Let\[\Lambda(k,m)=\limsup_{p\to \infty} r(k,m,p).\]Is it true that $\Lambda(k,2)$ is finite for all $k$? Is $\Lambda(k,3)$ finite for all odd $k$? How large are they?

\bigskip
\noindent\textbf{Remarks:}

Asked by Lehmer and Lehmer \cite{LeLe62}, who note that for example $\Lambda(2,2)=9$ - indeed, $9$ is always a quadratic residue, and if $10$ isn't then either $2$ or $5$ is, and hence at least one of $1,2$ or $4,5$ or $9,10$ is a consecutive pair of quadratic residues (and similarly there are infinitely many $p$ for which there are no consecutive quadratic residues below $9,10$).

A similar argument of Dunton \cite{Du65} proves $\Lambda(3,2)=77$, and Bierstedt and Mills \cite{BiMi63} proved $\Lambda(4,2)=1224$. Lehmer and Lehmer proved that $\Lambda(k,3)=\infty$ for all even $k$ and $\Lambda(k,4)=\infty$ for all $k\leq 1048909$.

Lehmer, Lehmer, and Mills \cite{LLM63} proved $\Lambda(5,2)=7888$ and $\Lambda(6,2)=202124$. Brillhart, Lehmer, and Lehmer \cite{BLL64} proved $\Lambda(7,2)=1649375$. Lehmer, Lehmer, Mills, and Selfridge \cite{LLMS62} proved that $\Lambda(3,3)=23532$. 

Graham \cite{Gr64g} proved that $\Lambda(k,l)=\infty$ for all $k\geq 2$ and $l\geq 4$.

Hildebrand \cite{Hi91} resolved the first question, proving that $\Lambda(k,2)$ is finite for all $k$: in other words, for any $k\geq 2$, if $p$ is sufficiently large then there exists a pair of consecutive $k$th power residues modulo $p$ in $[1,O_k(1)]$.

The remaining questions are to examine whether $\Lambda(k,3)$ is finite for all odd $k\geq 5$, and the growth rate of $\Lambda(k,2)$ and $\Lambda(k,3)$ as functions of $k$.

\bigskip
\noindent\textbf{References:}

[BLL64] Brillhart, John and Lehmer, D. H. and Lehmer, Emma, Bounds for pairs of consecutive seventh and higher power
residues. Math. Comp. (1964), 397--407.
 
 [BiMi63] Bierstedt, R. G. and Mills, W. H., On the bound for a pair of consecutive quartic residues of a
prime. Proc. Amer. Math. Soc. (1963), 628--632.
 
 [Du65] Dunton, M., Bounds for pairs of cubic residues. Proc. Amer. Math. Soc. (1965), 330--332.
 
 [Gr64g] Graham, R. L., On quadruples of consecutive {$k$}th power residues. Proc. Amer. Math. Soc. (1964), 196--197.
 
 [Hi91] Hildebrand, Adolf, On consecutive {$k$}th power residues. II. Michigan Math. J. (1991), 241-253.
 
 [LLM63] Lehmer, D. H. and Lehmer, Emma and Mills, W. H., Pairs of consecutive power residues. Canadian J. Math. (1963), 172--177.
 
 [LLMS62] Lehmer, D. H. and Lehmer, E. and Mills, W. H. and Selfridge,
J. L., Machine proof of a theorem on cubic residues. Math. Comp. (1962), 407--415.
 
 [LeLe62] Lehmer, D. H. and Lehmer, Emma, On runs of residues. Proc. Amer. Math. Soc. (1962), 102-106.

\bigskip
\noindent\small{Source: \url{https://www.erdosproblems.com/436}}
\end{document}
